\section{Discussion}

\textbf{Answer to the research question.} Yes, illicit and licit transactions
can be told apart, with an illicit F1 around 0.8 and PR-AUC around 0.8 on unseen
future time steps, and the transaction graph helps. But almost all of the graph's
value is captured by the simplest summary: the features of a transaction's direct
neighbours. Learning from the graph end-to-end (GNNs), adding structural
statistics or adding GNN embeddings does not beat a tree ensemble on those
features. The one graph idea that clearly paid off was respecting the
\emph{direction} of payments, which lifted a GNN by about 0.15 F1.

\textbf{What a practitioner should take away.}
\begin{enumerate}
  \item \emph{Use a tree ensemble on local + neighbour features.} Random Forest
        and XGBoost are tied, strong and cheap. Neither needs re-balancing.
  \item \emph{Report uncertainty over time, not over rows.} With 15 test steps,
        the 95\% interval of F1 is about 0.27 wide. Most published differences on
        this benchmark are smaller than that.
  \item \emph{Do not assume crime is anomalous.} Here it is the opposite:
        illicit transactions are typical, and the anomalies are licit.
  \item \emph{Plan for regime changes, not for average performance.} After the
        dark market shutdown every model fails, retraining recovers little, and
        the new illicit transactions look like licit ones. Monitoring the score
        distribution and fast labelling of new cases matter more than model
        choice.
  \item \emph{Check embeddings for leakage.} A hybrid that feeds a tree with
        embeddings of a GNN trained on the same labels is only safe while the GNN
        is weak; out-of-fold embeddings are required otherwise.
\end{enumerate}

\textbf{Limitations.}
(i) The features are anonymised, so we cannot say \emph{what} the model learned
in domain terms, only which group of features it uses.
(ii) After step 43 there are only \EdaIllicitPost{} illicit test transactions,
and step 46 has almost none; per-step scores there are noisy.
(iii) The ``typical'' Isolation Forest needs one bit of prior knowledge (the
direction), which we only know from labelled data; on another network the
direction may differ.
(iv) We tuned each GNN over a small grid (four configurations); a larger search
or temporal GNNs such as EvolveGCN~\cite{pareja2020} might narrow the gap to the
tree models, though they did not in the original study.
(v) Labels come from Elliptic's heuristics and are themselves imperfect; 77\% of
transactions have none.

\textbf{Future work.} Semi-supervised use of the 157,205 unlabelled
transactions; drift detectors that raise an alarm when the score distribution
shifts, so retraining is triggered before performance is lost; and the richer
Elliptic++ data~\cite{elmougy2023}, which adds wallet-level information.

\section{Conclusion}

We reproduced the Elliptic benchmark, then rebuilt its evaluation so that no
choice depends on the test period and every claim carries a confidence interval.
Under this protocol, tree ensembles on local and neighbour features remain the
best detectors (F1 \MainRFFone{}, PR-AUC \MainXGBPRAUC{}); the graph helps through
neighbour aggregation, and directed message passing is the most effective GNN
design. Two results contradict common assumptions: illicit transactions are
\emph{less} anomalous than licit ones, and a better graph model can make a hybrid
\emph{worse}. Finally, we showed why every model fails after the dark market
shutdown: the illicit population changed in almost every feature and came to
resemble licit traffic. Detecting financial crime on a blockchain is less a
modelling problem than a problem of keeping up with a moving target.
